% resumo em português
 % ajusta o espaçamento dos parágrafos do resumo
\begin{resumo}
    \SingleSpacing
     O presente trabalho apresenta o desenvolvimento do Avisa Aí, um aplicativo voltado ao registro e ao acompanhamento de denúncias relacionadas a problemas de infraestrutura, saneamento e demais questões que impactam diretamente a qualidade de vida dos munícipes de Itu. A definição do tema e o aprofundamento da problemática emergiram a partir da análise de dados epidemiológicos referentes ao aumento dos casos de dengue no município, assim como por meio de pesquisas quantitativas que identificaram os principais problemas apontados pela população ituana, revelando a necessidade de um canal direto, acessível e eficiente entre cidadãos e o poder público. Outro elemento norteador deste projeto foi o conceito de Cidades Inteligentes, que tem ganhado destaque no cenário mundial ao propor a utilização de tecnologias, dados e participação social para promover melhorias na gestão urbana e no bem-estar coletivo. A construção do aplicativo seguiu fundamentos da Engenharia de Software, com base nos modelos de requisitos, planejamento e processos descritos por Ian Sommerville. Para a implementação, adotaram-se princípios da Arquitetura Limpa, propostos por Robert C. Martin, visando modularidade, extensibilidade e testabilidade do sistema. A solução final consiste em uma aplicação mobile com interface responsiva e design intuitivo, destinada a facilitar o registro de denúncias por parte da população e proporcionar maior transparência nas ações administrativas. Dessa forma, o Avisa Aí se apresenta como uma ferramenta estratégica para fortalecer a participação cidadã, otimizar a comunicação entre população e prefeitura e contribuir para uma gestão pública mais eficaz, ágil e orientada por dados.
    \vspace{\onelineskip}

    \textbf{Palavras-chave}: Aplicativo mobile; denúncias de Infraestrutura; cidades inteligentes; arquitetura limpa..
\end{resumo}

% resumo em inglês
\begin{resumo}[Abstract]
    \begin{otherlanguage*}{english}
    \SingleSpacing
        This paper presents the development of Avisa Aí, an application aimed at registering and monitoring complaints related to infrastructure problems, sanitation, and other issues that directly impact the quality of life of Itu residents. The definition of the theme and the deepening of the problem emerged from the analysis of epidemiological data regarding the increase in dengue cases in the municipality, as well as through quantitative research that identified the main problems pointed out by the population of Itu, revealing the need for a direct, accessible, and efficient channel between citizens and the public authorities. Another guiding element of this project was the concept of Smart Cities, which has gained prominence on the world stage by proposing the use of technologies, data, and social participation to promote improvements in urban management and collective well-being. The construction of the application followed the fundamentals of Software Engineering, based on the requirements, planning, and process models described by Ian Sommerville. For implementation, principles of Clean Architecture, proposed by Robert C. Martin, were adopted, aiming at modularity, extensibility, and testability of the system. The final solution consists of a mobile application with a responsive interface and intuitive design, intended to facilitate the registration of complaints by the public and provide greater transparency in administrative actions. In this way, Avisa Aí presents itself as a strategic tool to strengthen citizen participation, optimize communication between the population and the city government, and contribute to a more effective, agile, and data-driven public administration.

        \vspace{\onelineskip}

        \noindent 
        \textbf{Keywords}: Mobile Application, Infrastructure Reporting, Smart Cities, Clean Architecture.
    \end{otherlanguage*}
\end{resumo}